% Proposed isolated A.93. All frozen cycle30 inputs remain unchanged.
\statement{A.93}{Proved}{The actual canonical and negative-hexagon equations in the three modules}
Let \(\alpha=(-1,f)\) be an actual nonidentity negative
\(\widehat{GT}_0\)-involution, with the canonical \(u\in F'\)
of A.85--A.89. Retain the original markings
\[
\begin{gathered}
 \alpha(x)=x^{-1},\quad \alpha(y)=f^{-1}y^{-1}f,
 \quad u=f\alpha(u),\quad z=(xy)^{-1},\\
 \iota(x)=y,\quad\iota(y)=x,\quad
 \iota(u)=u^{-1},\quad\iota(f)=f^{-1},\\
 v=\operatorname{Inn}(u^{-1})\iota,\qquad \tau=\alpha v.
\end{gathered}
\tag{A.93.1}
\]
Use the full characteristic-two procyclic coefficient ring
\(\Lambda=\mathbb F_2[[\widehat{\mathbb Z}]]\) in the three
cyclic-cover coordinates of A.92. Define
\[
\begin{gathered}
 U=[u]_x=[u]_y,\qquad W=[u]_z=[u(y,z)]_x,\\
 F_f=[f]_x=[f]_y,\qquad G_f=[f]_z=[f(y,z)]_x,
 \qquad j(P)=T^{-1}\overline P.
\end{gathered}
\tag{A.93.2}
\]
Then the actual canonical equation and the actual negative
hexagon imply
\[
             F_f=U+\overline U,\qquad G_f=0,
             \qquad j(W)=W.
\tag{A.93.3}
\]
The complete joint action on the direct sum of these three
kernel modules is
\[
\begin{aligned}
 \alpha(P_x,P_y,P_z)&=(\overline{P_x},\overline{P_y},j(P_z)),\\
 v(P_x,P_y,P_z)&=(P_y,P_x,j(P_z)),\\
 \tau(P_x,P_y,P_z)&=(\overline{P_y},\overline{P_x},P_z).
\end{aligned}
\tag{A.93.4}
\]
In these modules, the canonical equation is precisely
\[
 (U,U,W)=(F_f,F_f,0)+(\overline U,\overline U,j(W)).
\tag{A.93.5}
\]
Thus it supplies no additional relation between U and W at
this level once the displayed negative-hexagon consequences
are retained. This is not a redundancy claim for the full
nonlinear canonical equation.

There is also an independent exact source-word assertion.
For every odd n the joint map consisting of A.92's three
coordinates and the marked odd dihedral quotient
\[
 D_{2n}=\langle a,R\mid a^2=R^n=1,\ aRa=R^{-1}\rangle,
 \qquad x\mapsto a,\quad y\mapsto aR
\tag{A.93.6}
\]
has image \(I_n^3\times C_n\) on \(F'\). Every permitted
U,W and every rotation exponent k modulo n can be simultaneously
realized by one actual full-source duality word u. The same
holds at all odd heads simultaneously for compatible data.
Those realizing words are not asserted to be canonical
\(\widehat{GT}_0\) or \(\widehat{GT}\) parameters.

\begin{proof}
\emph{The original negative hexagon in its actual order.}\par
The full swap identities in (A.93.1), and the third-coordinate
action from A.92, give
\[
            \overline W=TW,\qquad \overline{G_f}=TG_f.
\]
On the x-cover, \(\alpha\) inverts the meridian and the
retained head. In characteristic two its module action is
coefficient bar. The conjugator f lies in the kernel and has
trivial conjugation action on its abelianization. Applying
\(u=f\alpha(u)\) gives
\[
                   F_f=U+\overline U,
                   \qquad \overline{F_f}=F_f.
\tag{A.93.7}
\]
The actual negative hexagon is the complete word
\[
        f(x,y)x^{-1}f(z,x)z^{-1}f(y,z)y^{-1}=1.
\tag{A.93.8}
\]
Reading its six factors in the x-cover, with the accumulated
heads retained, gives respectively
\[
                   F_f,quad1,quad\overline{F_f},
                        \quad1,quad TG_f,quad0.
\tag{A.93.9}
\]
For the third factor one uses the genuine substitution
\(x\mapsto x,y\mapsto z\), followed by full swap duality,
as in A.89; its action on this kernel module is coefficient
bar. The x-inverse contributes one, and the z-inverse contributes
one. Hence (A.93.8) gives
\(F_f+\overline{F_f}+TG_f=0\). Equation (A.93.7) and
invertibility of T imply \(G_f=0\). No positive multiplier d
has been inserted into this negative hexagon.

\emph{The joint Klein-four action.}\par
The y-cover calculation is the same as the x-cover calculation:
the meridian is inverted up to a kernel conjugator and the
retained x-head is inverted. On the z-cover, \(G_f=0\)
makes the image of f trivial. The original images there are
\(x=(0,T)\), \(y=(T^{-1},T^{-1})\), so
\[
 \alpha(x)=(0,T^{-1}),\quad
 \alpha(y)=(1,T),\quad \alpha(z)=(T^{-1},1).
\]
Thus its action on the z-kernel module is exactly j.
Conjugation by u has no effect on any of these kernel
abelianizations because u belongs to all three kernels.
Consequently v acts by the actual swap action of A.92, and
composition gives the stated \(\tau\)-action. These operators
indeed satisfy all Klein-four relations; they are obtained from
the actual group actions, rather than independently imposed.
Applying the canonical equation in all coordinates now gives
(A.93.5). Its first two entries are (A.93.7), and its last
entry is exactly the already known weighted duality of W.

\emph{The independent marked rotation coordinate.}\par
In (A.93.6), the commutator \([x,y]\) maps to \(R^2\).
As n is odd, \(F'\) therefore maps onto the full rotation
group \(C_n\). A.92 gives surjectivity onto \(I_n^3\).
Their joint image is the full direct product: if a source
word has image \((v,R^r)\), its nth power has image \((v,1)\),
since the first factor has exponent two. Squaring kills its
first coordinate, and multiplication by two is invertible
on \(C_n\). Hence the joint image contains both factors.

The actual swap inverts the rotation R. Given U,W,k, choose
the norm preimage \(P_z\) from A.92 and then an actual
\(w\in F'\) with joint coordinates
\[
                  (U,0,P_z;R^{k/2}).
\]
For \(u=w\iota(w)^{-1}\), the first three coordinates are
\((U,U,W)\), its rotation is \(R^k\), and
\(\iota(u)=u^{-1}\) is a full-group equality. The compact
finite-intersection argument over odd least common multiples
gives the simultaneous assertion, with w as the compact
variable. The trivial n=1 case imposes no extra condition.

For actual canonical data, the calculation above derives all
the displayed coefficient equations contributed by the
canonical involution condition and its negative hexagon in
these chosen modules. The separate source range theorem is
not a realization theorem for canonical data. In particular
the marked dihedral exponent k is not identified with the
different abelian-frame coordinate t of A.86. The full-GT
dyadic result of A.89, the remaining criterion of A.90, and
the nonlinear cusp equation retain their precise scopes.
\end{proof}
